% 04-visibility.tex -- the master reduction.  Formulas F9-F11.
% Input by mostlyharmless.tex.  No preamble here.

\section{The master reduction}
\label{sec:visibility}

\subsection{From the transform to a radial integral}
\label{sec:separation}

The complex visibility of a map is \eqref{eq:visibility-definition}.  Write the
disk in polar coordinates about North, $x = \rho\cos\varphi$ and
$y = \rho\sin\varphi$, and the baseline likewise, $u = q\cos\psi$ and
$v = q\sin\psi$, so that
\begin{equation}
  u x + v y = q\rho \cos(\varphi - \psi) ,
  \label{eq:phase-polar}
\end{equation}
with $q$ and $\psi$ from \eqref{eq:uv-polar}.  Both azimuths are measured from
North, so the difference $\varphi-\psi$ means what it always did; only the pairing
of symbol with axis has changed.  For a mode whose angular part is
$f(\rho)\,e^{im\varphi}$ the transform separates:
\begin{equation}
  V = \int_0^1 f(\rho)\,\rho
      \left[ \int_0^{2\pi} e^{im\varphi}\,
      e^{-2\pi i q\rho\cos(\varphi-\psi)} \, d\varphi \right] d\rho .
  \label{eq:separated}
\end{equation}
The inner integral is the Jacobi--Anger relation.  Because its integrand is
periodic, the uniform rule converges on it faster than any power of the step, and
certificate \vtag{17} checks both branches:
\begin{equation}
  \int_0^{2\pi} \cos(m\varphi)\, e^{-iz\cos(\varphi-\psi)} \, d\varphi
  = 2\pi(-i)^m J_m(z)\cos(m\psi) ,
  \label{eq:jacobi-anger}
\end{equation}
and the same statement with $\sin$ in place of $\cos$.  \emph{Numerical}, with a
residual at the roundoff level.

\subsection{The reduction}
\label{sec:master}

Substituting \eqref{eq:jacobi-anger} into \eqref{eq:separated} gives the form
everything else rests on.  For the complex mode $Y_{lm}$, whose angular part is
$N_{lm}P_{lm}(\cos\theta)\,e^{im\varphi}$,
\begin{equation}
  V_{lm}(u,v) = 2\pi(-i)^m e^{im\psi}\, N_{lm}
  \int_0^1 P_{lm}\!\left(\sqrt{1-\rho^2}\right) J_m(2\pi q\rho)\,\rho\,d\rho .
  \label{eq:master-reduction}
\end{equation}
Everything after this point is the evaluation of that radial integral: one
spherical Bessel function for the odd family, a Zernike sum for the even one.

Both factors in front of the integral are phases, so they combine.  Since
$-i = e^{-i\pi/2}$,
\begin{equation}
  (-i)^m e^{im\psi} = e^{im(\psi - \pi/2)} ,
  \label{eq:phase-packed}
\end{equation}
which is the compact way to remember the phase of a complex mode: it is $m$ times
the baseline azimuth, measured from a direction a right angle away from $+u$.

This cannot be carried over to the $\cos(m\psi)$ and $\sin(m\psi)$ factors of the
real basis, mentioned in section~\ref{sec:contribution}, and it is worth seeing why: each of those is a
pair of conjugate exponentials,
\[
  \cos(m\psi) = \tfrac12\left[ e^{im\psi} + e^{-im\psi} \right] ,
\]
so there is no single exponent to add the phase to.

The measure is the planar one of section 1.3, which is what makes
\eqref{eq:master-reduction} a half-range integral against a bare $\rho\,d\rho$.
The solid-angle alternative carries a factor $1/n$, and moves the square root
between the two families --- it swaps which of them is easy.

\subsection{The visibility of a single mode}
\label{sec:mode-visibility}

For given $l$ and $m \ge 0$, \eqref{eq:master-reduction} is a complex constant
times a radial integral that depends on $q$ alone.  That integral is the only object left to
evaluate, so it is given a letter of its own, $X_{lm}(q)$, deliberately not $R$: $R_n^m$ is the
Zernike radial polynomial and is needed in the same sections, and one letter cannot mean both.
\begin{equation}
  V_{lm}(u,v) = 2\pi(-i)^m N_{lm}\, e^{im\psi}\, X_{lm}(q) ,
  \label{eq:mode-visibility}
\end{equation}
\begin{equation}
  X_{lm}(q) = \int_0^1 P_{lm}\!\left(\sqrt{1-\rho^2}\right) J_m(2\pi q\rho)\,
  \rho\,d\rho .
  \label{eq:radial-integral}
\end{equation}
Written this way the reduction is complete and one object remains, $X_{lm}(q)$.
Sections 5 and 6 evaluate it for $l+m$ odd and even respectively, and that is the
whole of the remaining work.

\subsection{The contribution of a coefficient}
\label{sec:contribution}

The visibility is linear in the map, so it inherits the map's pairing of \eqref{eq:ylm-extension}.
One consequence is worth a sentence rather than an equation: at a negative order the visibility
carries \emph{no} extra sign, because the $(-1)^m$ of the basis is cancelled by the $(-1)^m$ that
$e^{im\pi}$ contributes when the baseline is reversed, as
\eqref{eq:pairing-visibility-parity} already showed.  That cancellation is the whole difference
between the map's inner sum and the visibility's, and it is why the stored half suffices on its
own.

The contribution of one stored coefficient is then not a new result but a substitution: take the
summand of \eqref{eq:pairing-visibility-parity} at order $m$ and put in the mode
\eqref{eq:mode-visibility}, which reduces both branches to a common form,
\[
  a_{lm}V_{lm} + (-1)^m\,\overline{a_{lm}V_{lm}}
  \;=\; 4\pi(-i)^m N_{lm} X_{lm}(q)\, \mathrm{Re}\!\left(a_{lm} e^{im\psi}\right) ,
  \qquad m > 0 ,
\]
because $(-i)^m$ is real for even $m$ and imaginary for odd $m$.  At $m = 0$ the contribution is
the single term $a_{l0}V_{l0}$.  The stored coefficients are complex; it is the map, and the
parameters of a fit, that are real.

Equations \eqref{eq:pairing-visibility-parity} and \eqref{eq:mode-visibility} are all that is
needed for the whole map.  The real basis requires no separate treatment: a branch is made of a
function and its conjugate, so its visibility follows from the one fact about conjugation used
throughout, and the results are collected with the rest of the published conventions in
section~\ref{sec:other-conventions}.

One could absorb \eqref{eq:phase-packed} into the azimuth for those branches by
measuring the baseline direction as a position angle, from North towards East, as is
conventional on the sky.  With $\PA_{uv} = \pi/2 - \psi$ the stray factor would
disappear.  It costs more than it saves, because the branches pair $\cos(m\varphi)$
with $\cos(m\psi)$: the body azimuth $\varphi$ is a longitude, not a position angle,
and re-expressing the baseline factor in the sky convention makes the cosine branch
contribute $\cos(m\PA_{uv} - m\pi/2)$, which is $\pm\sin(m\PA_{uv})$ for odd $m$ and
$\pm\cos(m\PA_{uv})$ for even $m$.  The branch names would alternate meaning with
$m$.  The two azimuths are measured the same way here on purpose.

The model in \code{make mapgui} draws exactly this: the real and imaginary parts of
$\Delta V$, in the observer's frame, on a common scale, with the origin at the centre of
each panel.  Both halves of \eqref{eq:pairing-visibility-parity} are then visible at once ---
the real part is even and the imaginary part odd under $(u,v) \to (-u,-v)$, since
$V(-u,-v) = \overline{V(u,v)}$ for a real map.

\subsection{The two parities}
\label{sec:parity}

Two reflections of the sphere, and their eigenvalues, govern almost everything that differs
between the mode families.  They are independent facts, and keeping them apart is what makes
the two closed forms and the several $(-1)^m$ factors legible.

\paragraph{Through the equator: the radial parity.}  Write $\rho = \sin\theta$, so
$\sqrt{1-\rho^2} = \cos\theta$, and reflect $\theta \to \pi - \theta$.  The angular factor
$P_{lm}(\cos\theta)$ is then either a polynomial in $\rho^2$, or that polynomial times a single
factor of $\sqrt{1-\rho^2}$, according to the parity of $l+m$:
\begin{equation}
  P_{lm}(\cos\theta) =
  \begin{cases}
    \text{polynomial in } \rho^2 , & l+m \text{ even} ,\\[2pt]
    \sqrt{1-\rho^2} \times \text{polynomial in } \rho^2 , & l+m \text{ odd} .
  \end{cases}
  \label{eq:parity-dichotomy}
\end{equation}
Equivalently, the radial profile is $\rho^m A(\rho^2)$ with $A$ a polynomial when $l+m$ is
even, and not otherwise.  That is why this parity decides the \emph{radial} answer.  For $l+m$
odd the half-range integral can be doubled into a full-range one, and it then collapses to a
single spherical Bessel function.  For $l+m$ even no such doubling exists and the result is a
finite sum of ordinary Bessel functions, obtained through a Zernike expansion.  Those are
sections~\ref{sec:odd} and~\ref{sec:even}.

\paragraph{About the axis: the angular parity.}  Reflect $\varphi \to \varphi + \pi$, a half
turn about the polar axis.  Only $m$ enters, through $e^{im\pi} = (-1)^m$, because the azimuthal
factor is the single exponential $e^{im\varphi}$ and not a sum of two.  That parity is what
governs the \emph{angular} behaviour, and it appears in four places: the phase $(-i)^m$ of the
mode, which is the statement $J_m(-z) = (-1)^mJ_m(z)$; the sign $(-1)^m$ in the pairing
$Y_{l,-m} = (-1)^m\overline{Y_{lm}}$; the switch between real and imaginary parts in
\eqref{eq:pairing-visibility-parity}; and the $\cos(m\psi)$ and $\sin(m\psi)$ factors of the real
branches.

The two are independent --- all four combinations of $l+m$ and $m$ occur --- which is why the
$(-1)^m$ of the angular parity turns up in arguments that have nothing to do with the
Bessel--Zernike split, and why it is worth naming them separately here.

Identity \eqref{eq:pairing-baseline-sign} is \textbf{certificate \vtag{23}}: it is checked
directly against \eqref{eq:mode-visibility} at the two baselines, with no coefficients and no
pairing involved, so the identity is isolated rather than inferred.  That isolation is the
point.  Hermitian symmetry of a real map, $V(-u,-v) = \overline{V(u,v)}$, also depends on it,
but only in combination with the pairing, so a sign error common to both would satisfy the
symmetry and still be wrong.  Certificate \vtag{23} would fail instead, and it would fail
loudly if the sense of the azimuth were ever changed again --- which it has been.

\subsection{The uniform disk}
\label{sec:uniform-disk}

A map with $a_{00} = 1$ and every other coefficient zero is the uniform brightness
case.  Only $l = m = 0$ survives \eqref{eq:master-reduction}, and its normalised
visibility must be the classical Airy visibility of a uniform disk,
\begin{equation}
  \frac{V(u,v)}{V(0,0)} = \frac{2J_1(x)}{x} , \qquad x = 2\pi q ,
  \label{eq:uniform-disk}
\end{equation}
for the unit disk of section 1.1.  The flux is $V(0,0) = \sqrt{\pi}/2$, since
$a_{00} = 1$ means unit mean intensity in the sense of
\eqref{eq:a00-normalisation}.

This is the anchor for the whole convention stack.  The harmonic normalisation,
the placement of $2\pi$ and the planar measure all enter \eqref{eq:uniform-disk},
and the classical result fixes them jointly.  Certificate \vtag{16} checks both
the flux and the normalised visibilities.  \emph{Numerical}, and exact in the
sense that matters: the map is one term, so no truncation is involved, and the
only error is the quadrature.

\subsection{What was checked}
\label{sec:visibility-register}

\begin{center}
\begin{tabular}{lll}
\hline
tag & claim & rung \\
\hline
\vtag{16} & uniform disk against $2J_1(x)/x$, and flux $\sqrt{\pi}/2$ & numerical, $l = 0$ \\
\vtag{17} & Jacobi--Anger factors \eqref{eq:jacobi-anger} & numerical, roundoff level \\
\hline
\end{tabular}
\end{center}

The radial integral in \eqref{eq:master-reduction} is evaluated by quadrature and
stays that way.  It is the reference against which the closed forms of the odd and
even families are checked, so its node count is stated in the code rather than
tuned silently, and a contract check compares it with an independent adaptive
quadrature.

One numerical detail is worth recording, because the contract check found it.  The
quadrature is done in $\theta$ after the substitution $\rho = \sin\theta$ of
section~\ref{sec:parity}, not in $\rho$.  Integrating in $\rho$ leaves the factor
$\sqrt{1-\rho^2}$ at the end point, where its derivative diverges; the integral
still converges, but only algebraically, and 64 Gauss--Legendre nodes give an
error near $10^{-6}$ for the odd-parity modes --- modes whose closed forms are the
subject of the next section.  In $\theta$ the integrand is analytic and the
convergence is limited only by the oscillation of the Bessel factor.
